\chapter{Metodologia}
\label{chap:metodologia}

\section{Modelos de Análise de Sobrevivência}

A abordagem proposta neste trabalho integra sinais de ECG Holter de 24 horas com
variáveis clínicas estruturadas em um modelo híbrido para a estimativa do tempo até
o evento. O pipeline é composto por três componentes principais: um encoder
convolucional profundo para extração de características dos sinais eletrocardiográficos,
uma rede de projeção para as variáveis clínicas e um cabeçote de predição baseado
em modelos de tempo discreto.

\subsection{Arquitetura do Encoder: ResNet-18 1D}

A extração de características dos sinais de ECG é realizada por uma versão
unidimensional da ResNet-18 \cite{he2016deep}, utilizada exclusivamente como
\textit{encoder}. O termo \textit{encoder} reflete o fato de que apenas o
corpo da rede é empregado --- ou seja, somente as camadas convolucionais
responsáveis pela extração de representações ---  sem a camada de
classificação final. Essa abordagem é recorrente na literatura quando se
deseja obter um vetor de características compacto para uso em tarefas
subsequentes, como a análise de sobrevivência proposta neste trabalho.

\subsubsection{Da ResNet-18 2D para a ResNet-18 1D}

A ResNet-18 original foi proposta para o processamento de imagens
bidimensionais \cite{he2016deep}, operando sobre tensores de entrada com
dimensões $[C, H, W]$, onde $C$ é o número de canais, $H$ a altura e $W$
a largura da imagem. A arquitetura é composta por 18 camadas com parâmetros
treináveis: uma convolução inicial com kernel $7{\times}7$, quatro grupos
de blocos residuais com 2 blocos cada, e uma camada totalmente conectada
(\textit{fully connected}) para classificação \cite{he2016deep}.

A adaptação para sinais unidimensionais consiste em substituir todas as
operações bidimensionais pelas suas equivalentes 1D, tornando a rede
adequada ao processamento de séries temporais. A entrada passa a ter
dimensões $[C, L]$, onde $C$ é o número de derivações do ECG --- 3 neste
trabalho --- e $L$ é o comprimento temporal do sinal. Todas as convoluções
2D são substituídas por convoluções 1D, e o \textit{batch normalization}
e o \textit{max pooling} são igualmente adaptados para operar sobre a
dimensão temporal. Essa estratégia é amplamente consolidada na literatura
de sinais biomédicos, com aplicações em ECG de derivação única
\cite{hannun2019cardiologist} e em ECG de 12 derivações
\cite{ribeiro2020automatic}.

\subsubsection{Estrutura das 17 Camadas do Encoder}

A ResNet-18 é composta originalmente por 17 camadas convolucionais e uma
camada totalmente conectada de classificação, totalizando 18 camadas com
parâmetros treináveis \cite{he2016deep}. Neste trabalho, a camada
totalmente conectada foi removida, e o encoder resultante --- com
\textbf{17 camadas convolucionais} --- é utilizado exclusivamente como
extrator de características, produzindo um vetor de 512 dimensões por
amostra ao final do processamento.

A arquitetura do encoder é organizada da seguinte forma:

\begin{enumerate}
	\item \textbf{Camada inicial} (\texttt{conv1}): convolução 1D com 64
	      filtros e kernel de tamanho 7, com \textit{stride} 2, seguida de
	      normalização em lote (\textit{batch normalization}) e ativação ReLU.
	      Uma camada de \textit{max pooling} com kernel 3 e \textit{stride} 2
	      reduz adicionalmente a resolução temporal, resultando em uma
	      representação inicial com 64 canais. Essa camada corresponde a
	      \textbf{1 camada convolucional}.

	\item \textbf{Grupo 1} (\texttt{layer1}, \texttt{conv2\_x}): dois
	      blocos residuais com 64 filtros e \textit{stride} 1, mantendo a
	      resolução temporal. Cada bloco contém 2 convoluções, totalizando
	      \textbf{4 camadas convolucionais}.

	\item \textbf{Grupo 2} (\texttt{layer2}, \texttt{conv3\_x}): dois
	      blocos residuais com 128 filtros e \textit{stride} 2, reduzindo a
	      resolução temporal à metade e dobrando a profundidade de representação.
	      Totalizando \textbf{4 camadas convolucionais}.

	\item \textbf{Grupo 3} (\texttt{layer3}, \texttt{conv4\_x}): dois
	      blocos residuais com 256 filtros e \textit{stride} 2, com nova redução
	      temporal e aprofundamento da representação. Totalizando \textbf{4
		      camadas convolucionais}.

	\item \textbf{Grupo 4} (\texttt{layer4}, \texttt{conv5\_x}): dois
	      blocos residuais com 512 filtros e \textit{stride} 2, produzindo a
	      representação de maior abstração semântica. Totalizando \textbf{4
		      camadas convolucionais}.
\end{enumerate}

Somando todas as contribuições: $1 + 4 + 4 + 4 + 4 = 17$ camadas
convolucionais, confirmando a estrutura do encoder utilizado.

Cada bloco residual (\texttt{BasicBlock1d}) é composto por duas convoluções
1D de kernel 3, cada uma seguida de normalização em lote e ativação ReLU,
com uma conexão de atalho (\textit{skip connection}) que soma a entrada do
bloco à sua saída. Quando há incompatibilidade de dimensões entre entrada e
saída --- o que ocorre nos grupos 2, 3 e 4 devido ao \textit{stride} 2 ---
uma projeção por convolução $1{\times}1$ é aplicada na conexão de atalho
para alinhar as dimensões. Ao final do encoder, um \textit{Global Average
	Pooling} reduz a dimensão temporal, produzindo um vetor fixo de 512
dimensões por amostra, independentemente do comprimento da série de entrada.

\subsubsection{Justificativa da Escolha da ResNet-18}

A escolha da ResNet-18 como backbone é motivada por dois fatores principais.
Primeiro, o conjunto de dados disponível é composto por aproximadamente 870
pacientes, um volume moderado para o treinamento de redes convolucionais
profundas. Arquiteturas mais profundas, como ResNet-34 ou ResNet-50, tendem
a demandar volumes maiores de dados para generalizar adequadamente,
tornando a ResNet-18 uma escolha mais equilibrada à escala do problema.
Segundo, a adaptação unidimensional da ResNet-18 é amplamente empregada
na literatura de sinais biomédicos \cite{hannun2019cardiologist,
	ribeiro2020automatic}, demonstrando capacidade suficiente para a extração
de padrões morfológicos relevantes em sinais de ECG com múltiplas
derivações.

Ao final do encoder, um \textit{Global Average Pooling} reduz a representação
temporal para um vetor fixo de 512 dimensões por amostra, independentemente do
comprimento da série de entrada.


\subsection{Módulo de Projeção e Fusão Híbrida}

Para equilibrar a contribuição dos sinais de ECG e das variáveis clínicas na
decisão final, ambas as modalidades são projetadas em espaços de mesma dimensão
antes da fusão. O vetor de 512 features extraídas pelo encoder é comprimido por
um \textit{bottleneck} com as camadas Linear(512, 64) $\to$ BatchNorm $\to$ ReLU
$\to$ Dropout $\to$ Linear(64, 16), resultando em um vetor de 16 dimensões.

As variáveis clínicas, compostas por 6 preditores numéricos, são processadas por
uma rede densa com as camadas Linear(6, 32) $\to$ BatchNorm $\to$ ReLU $\to$
Dropout $\to$ Linear(32, 16), produzindo igualmente um vetor de 16 dimensões.

Os dois vetores são então concatenados, formando uma representação híbrida de 32
dimensões que integra informação morfológica do ECG com o contexto clínico do
paciente.

\subsection{Cabeçote de Análise de Sobrevivência}

O vetor fundido de 32 dimensões é processado por um cabeçote de predição composto
pelas camadas Linear(32, 32) $\to$ BatchNorm $\to$ ReLU $\to$ Dropout $\to$
Linear(32, $K$), onde $K$ corresponde ao número de intervalos de tempo discretos
definidos pela transformação \textit{LabTrans}. A saída é um vetor de logits, cuja
interpretação varia conforme o modelo de sobrevivência selecionado, conforme
descrito nas subseções seguintes.


\subsection{Risco Logístico (Logistic-Hazard)}

O método \textit{Logistic-Hazard} parametriza diretamente os riscos discretos
$h(t_k \mid \mathbf{x})$ para cada intervalo de tempo $k$, otimizando a
verossimilhança de sobrevivência \cite{kvamme2019continuous, efron1988logistic}.
A função de risco discreta é obtida aplicando-se a função sigmoide sobre os
logits produzidos pelo cabeçote de predição, e a função de sobrevivência é
estimada pelo produto acumulado dos complementos dos riscos em cada intervalo.
Este método também é referenciado na literatura como Regressão Logística Parcial
\cite{biganzoli1998feedforward} e Nnet-sobrevivência (\textit{Nnet-survival})
\cite{gensheimer2019scalable}.

\subsection{Função de Massa de Probabilidade (PMF)}

O método PMF parametriza a função de massa de probabilidade do tempo até o evento,
otimizando diretamente a verossimilhança de sobrevivência \cite{kvamme2019continuous}.
A distribuição de probabilidade sobre os intervalos de tempo é obtida aplicando-se
a função \textit{softmax} sobre os logits da rede, e a função de sobrevivência é
estimada pela soma acumulada reversa dessa distribuição. Este método constitui a
base de abordagens como o DeepHit e o MTLR (\textit{Multi-Task Logistic Regression}).

\subsection{DeepHit}

O DeepHit é uma extensão do método PMF que incorpora, além da verossimilhança
de sobrevivência, um termo de perda por classificação baseado em \textit{ranking},
capaz de lidar com riscos concorrentes \cite{lee2018deephit}. A função de perda
total é definida como $\mathcal{L} = \mathcal{L}_1 + \alpha \cdot \mathcal{L}_2$,
onde $\mathcal{L}_1$ corresponde à log-verossimilhança negativa e $\mathcal{L}_2$
é um termo de \textit{ranking} que penaliza inversões na ordenação do risco entre
pares de pacientes com diferentes tempos de evento.

\subsection{Discretização do Tempo: Transformação LabTrans}

A análise de sobrevivência em tempo discreto requer que o tempo contínuo até o
evento seja mapeado em um conjunto finito de $K$ intervalos. Para isso, foi
implementada a transformação \texttt{DiscreteTimeTransform}, inspirada na
\textit{LabTrans} da biblioteca PyCox \cite{kvamme2019continuous}.

Dado um conjunto de tempos de sobrevivência $\{t_i\}$ e indicadores de evento
$\{d_i\}$, os pontos de corte $\mathbf{c} = (c_0, c_1, \ldots, c_K)$ são
determinados pelos quantis da distribuição dos tempos de evento observados
(i.e., $d_i = 1$), com $c_0 = 0$ e $c_K$ fixado acima do maior tempo
observado no conjunto. O índice de intervalo de cada observação é então
obtido por digitalização (\textit{digitize}), sendo clipado ao intervalo
$[0, K-1]$ para evitar extrapolações. O número de intervalos $K$ é um
hiperparâmetro do modelo, definido como \texttt{num\_durations}.

\subsection{Função de Perda por Tipo de Paciente}

Em dados de sobrevivência, as observações são heterogêneas: parte dos pacientes
apresentou o evento de interesse (\textit{uncensored}), enquanto os demais foram
censurados, ou seja, perderam o acompanhamento antes do desfecho. Cada modelo de
perda trata essas duas populações de forma distinta.

Para o \textit{Logistic-Hazard}, pacientes com evento contribuem com a
log-verossimilhança negativa do risco no instante $t_i$ ponderada pela
sobrevivência acumulada até $t_i - 1$. Pacientes censurados contribuem apenas
com a probabilidade de sobreviver até o instante de censura
\cite{kvamme2019continuous}. Formalmente:
%
\begin{equation}
	\mathcal{L}_{\text{LH}} = -\frac{1}{N} \sum_{i=1}^{N}
	\left[ d_i \left( \log h_{t_i} + \log S(t_i - 1) \right)
		+ (1 - d_i) \log S(t_i) \right]
\end{equation}
%
onde $h_{t_k} = \sigma(\phi_k)$ é o risco discreto no intervalo $k$, obtido
pela função sigmoide sobre o logit $\phi_k$, e $S(t_k) = \prod_{j \leq k}(1 - h_{t_j})$
é a função de sobrevivência estimada.

Para os métodos PMF e DeepHit, a distribuição de probabilidade sobre os intervalos
é obtida via \textit{softmax} sobre os logits, com um logit auxiliar nulo
concatenado ao final para representar sobrevivência além do horizonte de observação
\cite{kvamme2019continuous}. Pacientes com evento contribuem com $-\log f(t_i)$,
onde $f(t_k)$ é a probabilidade de evento no intervalo $k$, e pacientes censurados
contribuem com $-\log S(t_i^+)$, a probabilidade de sobreviver além do instante
de censura.

O DeepHit acrescenta um segundo termo de perda baseado em \textit{ranking}
\cite{lee2018deephit}:
%
\begin{equation}
	\mathcal{L}_{\text{DH}} = \mathcal{L}_1 + \alpha \cdot \mathcal{L}_2
\end{equation}
%
onde $\mathcal{L}_2$ penaliza pares de pacientes $(i, j)$ em que $t_i < t_j$
e $d_i = 1$, mas o modelo atribui risco cumulativo menor ao paciente $i$
do que ao paciente $j$ no tempo $t_i$. A penalidade é suavizada por uma
função exponencial com parâmetro $\sigma$:
%
\begin{equation}
	\mathcal{L}_2 = \sum_{i: d_i=1} \sum_{j} \mathbf{1}[t_i < t_j]
	\cdot \exp\!\left( -\frac{F(t_i \mid \mathbf{x}_i) - F(t_i \mid \mathbf{x}_j)}{\sigma} \right)
\end{equation}

\subsection{Treinamento: Acumulação de Gradiente e Precisão Mista}

O treinamento do modelo foi realizado com duas técnicas complementares para
viabilizar o uso de \textit{batches} efetivos maiores sem aumentar o consumo
de memória da GPU: acumulação de gradiente (\textit{gradient accumulation}) e
precisão mista automática (\textit{Automatic Mixed Precision} --- AMP)
\cite{micikevicius2018mixed}.

Na acumulação de gradiente, os gradientes são acumulados ao longo de
\texttt{accumulation\_steps} iterações antes de cada atualização dos pesos.
A função de perda de cada mini-\textit{batch} é dividida por
\texttt{accumulation\_steps} antes da retropropagação, de modo que a
magnitude dos gradientes acumulados seja equivalente à de um \textit{batch}
de tamanho \texttt{batch\_size} $\times$ \texttt{accumulation\_steps}. A
atualização dos pesos e o zeragem dos gradientes são realizadas apenas ao
final de cada ciclo de acumulação. Ao término do laço de treinamento, um
passo final é executado caso restem gradientes não aplicados.

A precisão mista é implementada com o \texttt{GradScaler} do PyTorch, que
escala dinamicamente a função de perda para evitar \textit{underflow} numérico
durante as operações em meia-precisão (\texttt{float16}), revertendo para
precisão simples (\texttt{float32}) nas etapas críticas de atualização. As
operações de \textit{forward} e cálculo da perda são executadas dentro de um
contexto \texttt{autocast}, enquanto a retropropagação e a atualização dos
pesos utilizam o escalonamento do \texttt{GradScaler}.

\subsection{Avaliação: C-Index e Agregação de Janelas}

A métrica principal de avaliação é o índice de concordância (\textit{C-index}),
que mede a capacidade discriminativa do modelo em ordenar corretamente pares de
pacientes pelo risco predito \cite{harrell1982evaluating}. O \textit{C-index}
varia entre 0,5 (aleatoriedade) e 1,0 (discriminação perfeita), sendo calculado
sobre os escores de risco derivados da função de sobrevivência estimada no último
intervalo de tempo.

Como os sinais de ECG Holter de 24 horas são segmentados em janelas para o
treinamento, cada paciente pode gerar múltiplas predições de risco. Para obter
um escore único por paciente na avaliação, foi implementada uma estratégia de
agregação com três variantes:

\begin{itemize}
	\item \textbf{Máximo} (\textit{max}): considera o maior escore de risco
	      entre todas as janelas do paciente, privilegiando a detecção de episódios
	      de alto risco isolados;

	\item \textbf{Média} (\textit{mean}): utiliza a média dos escores de risco,
	      representando o risco global ao longo do registro;

	\item \textbf{Percentil 90} (\textit{p90}): considera o percentil 90 da
	      distribuição de escores, oferecendo um compromisso entre a sensibilidade
	      a picos e a robustez a ruídos.
\end{itemize}

O escore agregado de cada paciente é então utilizado para o cálculo do
\textit{C-index} com a biblioteca \textit{lifelines}, tendo como referência
o tempo de sobrevivência e o indicador de evento de cada indivíduo.



Texto texto texto texto texto texto texto texto texto texto texto texto texto texto texto texto texto texto texto texto texto texto texto texto texto texto texto texto texto texto texto texto texto texto texto texto texto texto texto texto texto texto texto texto texto texto texto texto texto texto texto texto texto texto texto texto texto texto texto texto texto texto texto texto texto texto texto texto texto.

\section{Exemplo de alíneas}\label{sec:exemplo-de-algoritmos-e-figuras}

Texto texto texto texto texto texto texto texto texto texto texto texto texto texto texto texto texto texto texto texto texto texto texto texto texto texto texto texto texto texto texto texto texto texto texto texto texto texto texto texto texto texto texto texto texto texto texto texto texto texto texto texto texto texto texto texto texto texto texto texto texto texto texto texto texto texto texto texto texto.

%\begin{algorithm}[h!]
%	\SetSpacedAlgorithm
%	\caption{\label{exemplo-de-algoritmo}Como escrever algoritmos no \LaTeX2e}
%	\Entrada{o proprio texto}
%	\Saida{como escrever algoritmos com  Latex:}% \LaTeX2e }
%	\Inicio{
%		inicialização;
%		\Repita{fim do texto}{
%			leia o atual;
%			\Se{entendeu}{
%				vá para o proximo\;
%				próximo se torna o atual;}
%			\Senao{volte ao início da seção;}
%		}
%	}	
%\end{algorithm}

Texto texto texto texto texto texto texto texto texto texto texto.

%\begin{algorithm}[H]
%	\Entrada{o proprio texto}
%	\Saida{como escrever algoritmos com \LaTeX2e }
%	\Inicio{
%		inicialização\;
%		\Repita{fim do texto}{
%			leia o atual\;
%			\Se{entendeu}{
%				vá para o próximo\;
%				próximo se torna o atual\;}
%			\Senao{volte ao início da seção\;}
%		}
%	}
%	\caption{Exemplo de Algoritmo Versao 02}
%\end{algorithm}

%\begin{algorithm}
%	\begin{algorithmic}
%	\Entrada{o proprio texto}
%	\Saida{como escrever algoritmos com \LaTeX2e }	
%	\end{algorithmic}
%\end{algorithm}

Exemplo de alíneas com números:

\begin{alineascomnumero}
	\item Texto texto texto texto texto texto texto texto texto texto texto texto .
	\item Texto texto texto texto texto texto texto texto texto texto texto texto .
	\item Texto texto texto texto texto texto texto texto texto texto texto texto .
	\item Texto texto texto texto texto texto texto texto texto texto texto texto .
	\item Texto texto texto texto texto texto texto texto texto texto texto texto .
	\item Texto texto texto texto texto texto texto texto texto texto texto texto .
\end{alineascomnumero}

Texto texto texto texto texto texto texto texto texto texto texto texto texto texto texto texto texto texto texto texto texto texto texto texto texto texto texto texto texto texto texto texto texto texto texto texto texto texto texto texto texto texto texto texto texto texto texto texto texto texto texto texto texto texto texto texto texto texto texto texto texto texto texto texto texto texto texto texto texto.

Ou então figuras podem ser incorporadas de arquivos externos, como é o caso da \autoref{fig-grafico-1}. Se a figura que ser incluída se tratar de um diagrama, um gráfico ou uma ilustração que você mesmo produza, priorize o uso de imagens vetoriais no formato PDF. Com isso, o tamanho do arquivo final do trabalho será menor, e as imagens terão uma apresentação melhor, principalmente quando impressas, uma vez que imagens vetorias são perfeitamente escaláveis para qualquer dimensão. Nesse caso, se for utilizar o Microsoft Excel para produzir gráficos, ou o Microsoft Word para produzir ilustrações, exporte-os como PDF e os incorpore ao documento conforme o exemplo abaixo. No entanto, para manter a coerência no uso de software livre (já que você está usando LaTeX e abnTeX),  teste a ferramenta InkScape\index{InkScape}. ao CorelDraw\index{CorelDraw} ou ao Adobe Illustrator\index{Adobe! Illustrator}.  De todo modo, caso não seja possível  utilizar arquivos de imagens como PDF, utilize qualquer outro formato, como JPEG, GIF, BMP, etc.  Nesse caso, você pode tentar aprimorar as imagens incorporadas com o software livre \index{Gimp}Gimp. Ele é uma alternativa livre ao Adobe Photoshop\index{Adobe! Photoshop}.

\section{Usando Fórmulas Matemáticas}

Para escrever um símbolo matemático no texto, escreva símbolo entre cifrões, por exemplo, $\alpha$, $\beta$ e $\gamma$ são símbolo do alfabeto grego. Se você quiser inserir equações enumeradas, siga a estrutura de
\begin{equation}
	\label{eq:indices}
	k_{n+1} = n^2 + k_n^2 - k_{n-1}.
\end{equation}
Observe a pontuação, pois a equação faz parte da frase e do parágrafo. Como a equação faz parte da frase, não se utiliza o \textit{label} numérico \ref{eq:indices}.

Quando for citar a Equação \ref{eq:indices} novamente no texto, utiliza-se o \textit{label} numérico. Repare que a palavra ``Equação'' foi escrita com ``E'' maiúsculo.

Um exemplo de equações com frações é dado por
\begin{equation}
	\label{eq:fracao}
	\begin{aligned}
		x = a_0 + \cfrac{1}{a_1
			+ \cfrac{1}{a_2
				+ \cfrac{1}{a_3 + \cfrac{1}{a_4} } } }.
	\end{aligned}
\end{equation}

Texto texto texto texto texto texto texto texto texto texto texto texto texto texto texto texto texto texto texto texto texto texto texto texto texto texto texto texto texto texto texto texto texto texto texto texto texto texto texto texto texto texto texto texto texto texto texto texto texto texto texto texto texto texto texto texto texto texto texto texto texto texto texto texto texto texto texto texto texto
\begin{equation}
	\begin{aligned}
		k_{n+1} = n^2 + k_n^2 - k_{n-1}.
	\end{aligned}
\end{equation}

Texto texto texto texto texto texto texto texto texto texto texto texto texto texto texto texto texto texto texto texto texto texto texto texto texto texto texto texto texto texto texto texto texto texto texto texto texto texto texto texto texto texto texto texto texto texto texto texto texto texto texto texto texto texto texto texto texto texto texto texto texto texto texto texto texto texto texto texto texto
\begin{equation}
	\label{eq:trigo}
	\begin{aligned}
		\cos (2\theta) = \cos^2 \theta - \sin^2 \theta
	\end{aligned}.
\end{equation}

Texto texto texto texto texto texto texto texto texto texto texto texto texto texto texto texto texto texto texto texto texto texto texto texto texto texto texto texto texto texto texto texto texto texto texto texto texto texto texto texto texto texto texto texto texto texto texto texto texto texto texto texto texto texto texto texto texto texto texto texto texto texto texto texto texto texto texto texto texto
\begin{equation}
	\label{eq:matriz}
	\begin{aligned}
		A_{m,n} =
		\begin{pmatrix}
			a_{1,1} & a_{1,2} & \cdots & a_{1,n} \\
			a_{2,1} & a_{2,2} & \cdots & a_{2,n} \\
			\vdots  & \vdots  & \ddots & \vdots  \\
			a_{m,1} & a_{m,2} & \cdots & a_{m,n}
		\end{pmatrix}
	\end{aligned}.
\end{equation}

Texto texto texto texto texto texto texto texto texto texto texto texto texto texto texto texto texto texto texto texto texto texto texto texto texto texto texto texto texto texto texto texto texto texto texto texto texto texto texto texto texto texto texto texto texto texto texto texto texto texto texto texto texto texto texto texto texto texto texto texto texto texto texto texto texto texto texto texto texto
\begin{equation}
	\label{eq:sistema}
	\begin{aligned}
		f(n) = \left\{
		\begin{array}{l l}
			n/2      & \quad \text{if $n$ is even} \\
			-(n+1)/2 & \quad \text{if $n$ is odd}
		\end{array} \right.
	\end{aligned}.
\end{equation}
Texto texto texto texto texto texto texto texto texto texto texto texto texto texto texto texto texto texto texto texto texto texto texto texto texto texto texto texto texto texto texto texto texto texto texto texto texto texto texto texto texto texto texto texto texto texto texto texto texto texto texto texto texto texto texto texto texto texto texto texto texto texto texto texto texto texto texto texto texto

%\section{Usando Algoritmos}

%\begin{algorithm}[h!]
%	\SetSpacedAlgorithm
%	\caption{\label{alg:algoritmo_de_colonica_de_formigas}Algoritmo de Otimização por Colônia de Formiga}
%	\Entrada{Entrada do Algoritmo}
%	\Saida{Saida do Algoritmo}
%	\Inicio{
%		Atribua os valores dos parâmetros\;
%		Inicialize as trilhas de feromônios\;
%		\Enqto{não atingir o critério de parada}{
%			\Para{cada formiga}{
%				Construa as Soluções\;
%			}
%			Aplique Busca Local (Opcional)\;
%			Atualize o Feromônio\;
%		}	
%	}		
%\end{algorithm}

\section{Usando Código-fonte}

Um exemplo de código-fonte, ou código de programação encontra-se no Apendice \ref{ap:A}


\section{Usando Teoremas, Proposições, etc}

Texto texto texto texto texto texto texto texto texto texto texto texto texto texto texto texto texto texto texto texto texto texto texto texto texto.

\begin{teo}[Pitágoras]
	Em todo triângulo retângulo o quadrado do comprimento da
	hipotenusa é igual a soma dos quadrados dos comprimentos dos catetos. Usando o Apêndice \ref{ap:C}
\end{teo}


Texto texto texto texto texto texto texto texto texto texto texto texto texto texto texto.

\begin{teo}[Fermat]
	Não existem inteiros $n > 2$, e $x, y, z$ tais que $x^n + y^n = z$
\end{teo}

Texto texto texto texto texto texto texto texto texto texto texto texto texto texto texto.

\begin{prop}
	Para demonstrar o Teorema de Pitágoras...
\end{prop}

Texto texto texto texto texto texto texto texto texto texto texto texto texto texto texto.

\begin{exem}
	Este é um exemplo do uso do ambiente exem definido acima.
\end{exem}

Texto texto texto texto texto texto texto texto texto texto texto texto texto texto texto.


\begin{xdefinicao}
	Definimos o produto de ...
\end{xdefinicao}

Texto texto texto texto texto texto texto texto texto texto texto texto texto texto texto.

\section{Usando Questões}

Um exemplo de questionário encontra-se no Apêndice \ref{ap:B}.

%Movido para o Apêndice

